
\section{Analysis}
\label{sec:analysis}

We will first state a formal version of Theorem~\ref{thm:regret}.
Recall from the main text where we stated that most evaluations at $\zhf$ are inside
the following set $\Xcalrho$.
\begin{align*}
\Xcalrho = \{x\in\Xcal: \funcopt - \func(x) \leq 2\rho\sqrtks\inffunc(p)\}.
\end{align*}
This is not entirely accurate as it hides a dilation that arises due to a covering
argument in our proofs.
Precisely, we will show that after $n$ queries at any fidelity, 
\bocas will use most of the $\zhf$ evaluations 
in $\Xcalrhon$ defined below using $\Xcalrho$.
\begin{align*}
\Xcalrhon = \big\{\,x\in\Xcal: {\rm B}_2\big(x, \sqrt{d}/n^{\alpha/2d}\big)
\cap \Xcalrhon \neq \emptyset \big\}
\numberthis \label{eqn:Xcalrhon}
\end{align*}
Here ${\rm B}_2(x, \epsilon)$ is an $L_2$ ball of radius $\epsilon$ centred at $x$.
$\Xcalrhon$ is a dilation of $\Xcalrho$ by $\sqrt{d}/n^{\alpha/2d}$.
Notice that for all $\alpha > 0$, as $n\rightarrow\infty$,
$\Xcalrhon$ approaches $\Xcalrho$ at a polynomial rate.
We now state our main theorem below.

\vspace{0.1in}

\begin{theorem}
\label{thm:main}
Let $\Zcal=[0,1]^p$ and $\Xcal=[0,1]^d$. Let $\gunc\sim\GP(\zero,\kernel)$
where $\kernel$ is of the form~\eqref{eqn:mfkernel}. Let $\phix$ satisfy 
Assumption~\ref{asm:kernelAssumption} with some constants $a,b>0$.
% and satisfy assumptions \emph{\Gtwo}, \emph{\Gthree}.
% Let $\kernel$ satisfy
% Assumption~\ref{asm:kernelAssumption} with some constants $a,b$.
Pick $\delta\in(0,1)$ and run \emph{\bocas} with
\[
\betat =
2\log\left(\frac{\pi^2t^2}{2\delta}\right) + 4d\log(t) + 
\max\left\{\,0\,,\, 
2d\log\left(brd\log\bigg(\frac{6ad}{\delta}\bigg) \right)\right\}.
\]
Then, for all $\alpha\in(0,1)$ 
there exists $\rho,\capital_0$ such that with probability at least $1-\delta$ we have
for all $\capital \geq \capital_0$,
\begin{align*}
S(\capital)\;&\leq\;
  \sqrt{\frac{2C_1\beta_{2\ncapital} \IGnn{2\ncapital}(\Xcalrhon)}{\ncapital} }
 \;+\;
  \sqrt{\frac{2C_1\beta_{2\ncapital} \IGnn{2\ncapital^\alpha}(\Xcal)}{\ncapital^{2-\alpha}} }
  \;+\; \frac{\pi^2}{6\ncapital}.
\end{align*}
Here $C_1 = 8/\log(1+\eta^2)$ is a constant and $\ncapital=\floor{\capital/\cost(\zhf)}$.
$\rho$ satisfies %depends only on $\alpha$ and satisfies
$\rho > \rho_0 = \max\{2, 1 + \sqrt{(1+2/\alpha)/(1+d)}\}$.
% $\COSTnought$ depends on several problem dependent quantities including
% the volumes of the sets $\Htcalm$ and consequently on $\rho$.
\end{theorem}
In addition to the dilation,
Theorem~\ref{thm:regret} in the main text also suppresses the constants and $\polylog$ terms.
The next three subsections are devoted to proving the above theorem.
In Section~\ref{sec:setup} we describe some discretisations for $\Zcal$ and $\Xcal$ which
we will use in our proofs.
Section~\ref{sec:techLemmas} gives some lemmas we will need and
Section~\ref{sec:mainproof} gives the proof.

% $\Xcalrhon$ arises due to a covering argument in our proofs.



\subsection{Set Up \& Notation}
\label{sec:setup}

\textbf{Notation:}
Let $U\subset\Zcal\times\Xcal$.
$\Tn(U)$ will denote the number of queries by \bocas at points $(z,x)\in U$ within
$n$ time steps.
When $A\subset\Zcal$ and $B\subset\Xcal$, we will overload notation to denote
$\Tn(A,B) = \Tn(A\times B)$.
% denotes the number of queries by \bocas
% at points $(z,x)\in W\times A$ within time $n$.
For $z\in\Zcal$, $\greaterz$  will denote the fidelities which are more expensive than $z$,
i.e. $\greaterz = \{z'\in\Zcal: \cost(z') > \cost(z)\}$.

We will require a fairly delicate set up before we can prove Theorem~\ref{thm:main}.
Let $\alpha >0$.
All sets described in the rest of this subsection are defined with respect to $\alpha$.
First define
\[
\Hcalntilde = \{(z,x)\in\Zcal\times\Xcal: \funcopt - \func(x) <
2\rho\betanh\sqrtks\inffunc(z) \},
\]
where recall from~\eqref{eqn:inffunc}, $\inffunc(z) = \sqrt{1-\phiz^2(\|z-\zhf\|)}$ is the
information gap function.
We next define $\Hcalnprime$ to be an $L_2$ dilation of $\Hcalntilde$ in the $\Xcal$
space, i.e.
\[
\Hcalnprime = 
\{ (z,x) \in \Zcal\times\Xcal:
{\rm B}_2\big(x, \sqrt{d}/n^{\alpha/2d}\big) \cup
\Hcalntilde \neq \emptyset
 \}.
\]
Finally, we define $\Hcaln$ to be the intersection of $\Hcalnprime$ with all fidelities
satisfying the third condition in~\eqref{eqn:candfidelst}.
That is,
\begin{align*}
\Hcaln = \Hcalnprime \cap \Big\{(z,x)\in\Zcal\times\Xcal: \inffunc(z) >
\inffunc(\sqrt{p})/\betanh \Big\}.
\numberthis \label{eqn:Hcaln}
\end{align*}
In our proof we will use the second condition in~\eqref{eqn:candfidelst} to control the
number of queries in $\Hcaln$.

To control the number of queries outside $\Hcaln$ we first introduce 
a $\frac{\sqrt{d}}{2n^{\frac{\alpha}{2d}}}$-covering of
the space $\Xcal$ of size $n^{\alpha/2}$.
If $\Xcal=[0,1]^d$, a sufficient covering would be an equally spaced grid having
$n^{\frac{\alpha}{2d}}$ points per side.
Let $\{\ain\}_{i=1}^{\nalphabt}$ be the points in the covering.
$\Ain\subset\Xcal$ to be the points in $\Xcal$ which are closest to $\ain$ in $\Xcal$.
Therefore $\Fn = \{\Ain\}_{i=1}^{\nalphabt}$ is a  partition of $\Xcal$.
% We now define $\Fcaln$ to be the union of all $\Ain$ which are outside of $\Hcalntilde$.

Now define $\Qt:2^\Xcal\rightarrow 2^\Zcal$ to be the following function
which maps subsets of $\Xcal$ to subsets of $\Zcal$.
\begin{align*}
\Qt(A) = \Big\{z\in\Zcal: \;\forall\,x\in A, \quad
      \funcopt - \func(x) \geq 2\rho\betath\sqrtks\inffunc(z) \Big\}.
\numberthis\label{eqn:Qt}
\end{align*}
That is, $\Qt$ maps $A\subset\Xcal$ to fidelities where the information gap $\inffunc$
is smaller than $(\funcopt - \func(x))/(2\rho\betath)$ for all $x\in A$.
Next we define $\thetat:2^\Xcal\rightarrow \Zcal$, to be the cheapest fidelity
in $\Qt(A)$ for a subset $A\in\Xcal$. 
\begin{align*}
\thetat(A)\,=\, \arginf_{z\in\Qt(A)} \cost(z).
\numberthis\label{eqn:thetat}
\end{align*}
We will see that \bocas will not query inside an $\Ain\in\Fn$ at fidelities larger
than $\thetat(\Ain)$ too many times (see Lemma~\ref{lem:gtrthetatbound}).
That is, $\Tnzzxx{\greaterzz{\thetan(\Ain)}}{\Ain}$ will be small.
We now define $\Fcaln$ as follows,
\begin{align*}
\Fcaln = \bigcup_{\Ain\subset\Xcal\setminus\Xcalrhon} 
\greaterzz{\thetan(\Ain)} \times \Ain.
\numberthis \label{eqn:Fcaln}
\end{align*}
That is, we first choose $\Ain$'s that are completely outside $\Xcalrhon$
and take their cross product with fidelities more expensive than $\thetat(\Ain)$.
By design of the above sets, and using the third condition in~\eqref{eqn:candfidelst}
we can bound the total number of queries as follows,
% $\Hcaln\cup\Fcaln\cup(\{\zhf\}\times\Xcalrhon) 
% = \Zcal\times\Xcal$. Therefore,
\begin{align*}
n = \Tn(\Zcal,\Xcal) \;\leq\;
  \Tn(\{\zhf\},\Xcalrhon)
  + \Tn(\Fcaln) + \Tn(\Hcaln)
\label{eqn:Tnbreakdown}
\end{align*}
We will show that the last two terms on the right hand side are small for \bocas
and consequently, the first term will be large.
But first, we establish a series of technical results which will be useful in
proving theorem~\ref{thm:main}.



\subsection{Some Technical Lemmas}
\label{sec:techLemmas}

The first lemma proves that the UCB $\utilt$ in~\eqref{eqn:mfucb}
upper bounds $\func(\xt)$ on all the domain points $\{\xt\}_{t\geq 1}$ chosen for
evaluation.
\insertprespacing
% 
\begin{lemma}
\label{lem:ucblemma}
Let $\betat > 2\log(\pi^2t^2/2\delta)$. Then, with probability $>1-\delta/3$, we have
\[
\forall\,t\geq 1,\quad
|\func(\xt) - \mutmo(\xt)|\,\leq\,\betath\sigmatmo(\xt).
\]
\end{lemma}
\textbf{Proof:}
This is a straightforward argument using Lemma~\ref{lem:gaussConcentration}
and the union bound.
At $t\geq 1$,
\begingroup
\allowdisplaybreaks
\begin{align*}
\PP\Big(|\func(x)-\mutmo(x)|>\betath\sigmatmo(x)\Big) &=
\EE\left[\EE\left[ |\func(x)-\mutmo(x)|>\betath\sigmatmo(x) 
  \;\Big|\; \Dcal_{t-1}  \right]\right] \\
&=\EE\left[ \PP_{Z\sim\Ncal(0,1)}\left(|Z|>\betath\right) \right] 
\leq\; \exp\Big(\frac{-\betat}{2}\Big)
\;=\; \frac{2\delta}{\pi^2t^2}.
\end{align*}
\endgroup
In the first step we have conditioned w.r.t $\Dcal_{t-1} = \{(z_i,x_i,y_i)\}_{i=1}^{t-1}$
which allows us to use Lemma~\ref{lem:gaussConcentration}
as $\func(x)|\Dcal_{t-1} \sim\Ncal(\mutmo(x),\sigmasqtmo(x))$.
The statement follows via a union bound over all 
$t\geq 0$ and the fact that $\sum_{t}t^{-2} = \pi^2/6$.
\hfill\BlackBox


Next we show that the GP sample paths are well behaved and that $\utilt(x)$ upper
bounds  $\func(x)$ on a sufficiently dense subset at each time step.
For this we use the following lemma.

\insertprespacing
\begin{lemma}
\label{lem:smoothnesslemma}
Let $\betat$ be as given in Theorem~\ref{thm:main}.
Then for all $t$, there exists a discretisation $G_t$ of $\Xcal$ of size
$(t^2brd\sqrt{6ad/\delta})^d$ such that the following hold.
\vspace{-0.1in}
\begin{itemize}
\item
Let $[x]$ be the closest point to $x\in\Xcal$ in the discretisation.
With probability $>1-\delta/6$, we have
\[
\forall\;t\geq 1,\quad
\forall\;x\in \Xcal,\quad
|\func(x) - \func([x]_t)| \leq 1/t^2. 
\]
\vspace{-0.25in}
\item
With probability $>1-\delta/3$,
for all $t\geq 1$ and for all $a\in G_t$,
$|\func(a) -\mutmo(a)| \leq \betath\sigmatmo(a)$.
\end{itemize}
\end{lemma}
\insertpostspacing
\textbf{Proof:}
The first part of the proof, which we skip here, uses the regularity condition
for $\phix$ in Assumption~\ref{asm:kernelAssumption} and mimics the argument in Lemmas
5.6, 5.7 of~\citet{srinivas10gpbandits}.
The second part mimics the proof of Lemma~\ref{lem:ucblemma} and uses the fact that
$\betat > 2\log(|G_t|\pi^2t^2/2\delta)$.
\hfill\BlackBox


The discretisation in the above lemma is different to the coverings introduced
in Section~\ref{sec:setup}.
The next lemma is about the information gap function in~\eqref{eqn:inffunc}.

\begin{lemma}
\label{lem:inffuncbound}
Let $g\sim\GP(0,\kernel)$, $g:\Zcal\times\Xcal\rightarrow\RR$ and $\kernel$ is
of the form~\eqref{eqn:mfkernel}.
Suppose we have $s$ observations from $g$.
Let $z\in\Zcal$ and $x\in\Xcal$.
Then $\tautmo(z,x) < \alpha$ implies $\sigmatmo(x) < \alpha + \sqrtks\inffunc(z)$.
\end{lemma}
\textbf{Proof:}
The proof uses the observation that for radial kernels, the maximum difference
between the variances at two points $u_1$ and $u_2$ occurs when all $s$ observations
are at $u_2$ or vice versa.
Now we use $u_1 = (z,x)$ and $u_2 = (\zhf,x)$ and apply Lemma~\ref{lem:varbound}
to obtain $\tausqtmo(\zhf,x) \leq \kernelscale(1-\phiz(\|\zhf-z\|))^2 +
\frac{\eta^2/s}{1 + \frac{\eta^2}{s\kernelscale}}$.
However, As $\tausqtmo(z,x) = \frac{\eta^2/s}{1 + \frac{\eta^2}{s\kernelscale}}$
when all observations are at $(z,x)$ and noting that
$\sigmasqtmo(x) = \tausqtmo(\zhf,x) $, we have
$\sigmasqtmo(\zhf,x) \leq \kernelscale(1-\phiz(\|\zhf-z\|))^2 + \tausqtmo(z,x)$.
Since the above situation characterised the maximum difference between
$\sigmasqtmo(x)$ and $\tausqtmo(z,x)$, this inequality is valid for any general
observation set.
The proof is completed using the elementary inequality $a^2 + b^2 \leq (a+b)^2$ for $a,b>0$.
\hfill\BlackBox


We are now ready to prove Theorem~\ref{thm:main}.
The plan of attack is as follows.
We will analyse \bocas after $n$ time steps and bound the number of plays at fidelities
$z\neq\zhf$ and outside $\Xcalrhon$ at $\zhf$.
Then we will show that for sufficiently large $\capital$, the number of \emph{random}
plays $N$ is bounded by $2\ncapital$ with high probability.
Finally we usee techniques from~\citet{srinivas10gpbandits}, specifically the
maximum information gain, to control the simple regret.
However, unlike them we will obtain a tighter bound as we can control the regret due
to the sets $\Xcalrhon$ and $\Xcal\setminus\Xcalrhon$ separately.

% \vspace{0.1in}
% \subsection{Set Up}

\subsection{\textbf{Proof of Theorem~\ref{thm:main}}}
\label{sec:mainproof}

Let $\alpha >0$ be given.
We invoke the sets $\Xcalrhon, \Hcaln,\Fcaln$ in equations~\eqref{eqn:Xcalrhon},
~\eqref{eqn:Hcaln},~\eqref{eqn:Fcaln}
for the given $\alpha$.
The following lemma establishes that for any $A\subset\Xcal$, we will not query
inside $A$ at fidelities larger than $\thetat(A)$~\eqref{eqn:thetat}
too many times.
The proof is given in Section~\ref{sec:gtrthetatbound}.
\begin{lemma}
\label{lem:gtrthetatbound}
Let $A\subset\Xcal$ which does not contain the optimum.
Let $\rho,\betat$ be as given in Theorem~\ref{thm:main}.
Then for all $u>\max\{3, (2(\rho-\rho_0)\eta)^{-2/3}\}$, we have
\[
\PP\Big(\Tnzzxx{\greaterzz{\thetat(A)}}{A} \,>\, u \Big)
\;\leq\; \frac{\delta}{\pi^2}\frac{1}{u^{1+4/\alpha}}
\]
\end{lemma}
To bound $T(\Fcaln)$, we will apply Lemma~\ref{lem:gtrthetatbound} with $u=n^{\alpha/2}$
on all $\Ain\in \Fn$ satisfying
$\Ain\subset\Xcal\setminus\Xcalrhon$.
Since $\Xcalrho\subset\Xcalrhon$, $\Ain$ does not contain the optimum.
As $\Fcaln$ is the union of such sets~\eqref{eqn:Fcaln}, we have for all
$n$ (larger than a constant),
\begin{align*}
\PP(T(\Fcaln) > n^{\alpha}) \;&\leq\;\;
\PP\Big( \exists \Ain\subset\Xcal\setminus\Xcalrhon, \;\;
\Tnzzxx{\greaterzz{\thetat(\Ain)}}{\Ain} \,>\,  n^{\alpha/2} \Big) \\
  &\leq
  \sum_{\substack{\Ain\in \Fn\\ \Ain\subset\Xcal\setminus\Xcalrhon}} \PP\Big(
    \Tnzzxx{\greaterzz{\thetat(\Ain)}}{\Ain} \,>\, n^{\alpha/2} \Big) 
  \leq |\Fn|\frac{\delta}{\pi^2} \frac{1}{n^{\alpha/2 + 2}} 
    \leq \frac{\delta}{\pi^2}\frac{1}{n^2}
\end{align*}
Now applying the union bound over all $n$, we get
$\PP(\forall\,n\geq 1,\;T(\Fcaln) > n^{\alpha}) \,\leq\, \delta/6$.

Now we will bound the number of plays in $\Hcaln$ using the second condition
in~\eqref{eqn:candfidelst}.
We begin with the following Lemma.
The proof mimics the argument in Lemma 11 of~\citet{kandasamy16mfbo} who prove
 a similar result for GPs defined on just the domain, i.e. $\func\sim\GP(\zero,\kernel)$
where $\func:\Xcal\rightarrow\RR$.

\begin{lemma}
\label{lem:mfvarbound}
Let $A\subset\Zcal\times\Xcal$ and the $L_2$ diameter of $A$ in $\Xcal$ be
$\Dx$ and that in $\Zcal$ be $\Dz$.
Suppose we have $n$ evaluations of $g$ of which $s$ are in $A$. 
Then for any $(z,x)\in A$, the posterior variance $\tau'^2$ satisfies, 
\[
\tau'^2(z,x) \leq \kernelscale(1 - \phiz^2(\Dz) \phix^2(\Dx)) + \frac{\eta^2}{s}.
\]
\end{lemma}
Let $\costratio = \cost_{min}/\cost(\zhf)$ where $\cost_{min} = \min_{z\in\Zcal}\cost(z)$.
% To show that we will not query in a certain region at time $n$, it is sufficient to prove
If the maximum posterior variance in a certain region is smaller than $\gamma(z)$,
then we will not query within that region by the second condition
in~\eqref{eqn:candfidelst}.
Further by the third condition, since we will only query at fidelities satisfying
 $\inffunc(z) > \inffunc(\sqrt{p})/\betanh$,
it is sufficient to show that
the posterior variance is bounded by
$\kernelscale \inffunc(\sqrt{p})^2 \costratio^{2q}/\beta_n$ at time $n$ 
to prove that we will not query again in that region.
% This uses the fact that $\inffunc(z) > \inffunc(\sqrt{p})/\betanh$ by the third condition.
For this we can construct a covering of $\Hcaln$ such that 
$1 - \phiz^2(\Dz) \phix^2(\Dx) < \frac{1}{2}\inffunc(\sqrt{p})^2\costratio^{2q}/\betan$.
For any $A\subset\Zcal\times\Xcal$, the covering number, which we denote $\Omega_n(A)$
of this construction
will typically be polylogarithmic in $n$ (See Remark~\ref{rem:construction} below).
Now if there are 
$\frac{2\betan\eta^2}{\costratio^{2q}\inffunc^2(\sqrt{p})\kernelscale} + 1$
queries inside a ball in this covering,
the posterior variance, by Lemma~\ref{lem:mfvarbound} will be smaller than
$\kernelscale \inffunc(\sqrt{p})^2 \costratio^{2q}/\beta_n$. Therefore, 
we will not query any further inside this ball.
Hence, the total number of queries in $\Hcaln$ is
$\Tn(\Hcaln) \leq C_2 \Omega_n(\Hcaln) \frac{\betan}{\costratio^{2q}}
\leq C_3 \vol(\Hcaln) \frac{\polylog(n)}{\poly(\costratio)}$
for appropriate constants $C_2,C_3$.
(Also see Remark~\ref{rem:qchoice}).

Next, we will argue that the number of queries for sufficiently large $\capital$,
is bounded by $\ncapital/2$ where, recall $\ncapital = \floor{\capital/\cost(\zhf)}$.
This simply follows from the bounds we have for $\Tn(\Fcaln)$ and $\Tn(\Hcaln)$.
\[
\Tn(\Zcal\setminus\{\zhf\},\Xcal)
\leq \Tn(\Fcaln) + \Tn(\Hcaln) \leq n^\alpha + \bigO(\polylog(n)).
\]
Since the right hand side is sub-linear in $n$, we can find $n_0$ such that
for all $n_0$, $n/2$ is larger than the right hand side.
Therefore for all $n\geq n_0$, $\Tn(\{\zhf\},\Xcal) > n/2$.
Since our bounds hold with probability $>1-\delta$ for all $n$ we can invert the
above inequality to bound $N$, the random number of queries after capital $\capital$.
We have $N\leq 2\capital/\cost(\zhf)$. We only need to make sure that $N\geq n_0$ which
can be guaranteed if $\capital > \capital_0 = n_0\cost(\zhf)$.

The final step of the proof is to bound the simple regret after $n$ time steps
in \boca. This uses techniques that are now standard in GP bandit optimisation, so we
only provide an outline. We begin with the following Lemma.
\begin{lemma}
\label{lem:IGBoundLemma}
Assume that we have queried $g$ at $n$ points, $(z_t,x_t)_{t=1}^n$ of
which $s$ points are in $\{\zhf\}\times A$ for any $A\subset\Xcal$.
Let $\sigmatmo$ denote the posterior variance of $\func$ at time $t$,
i.e. after $t-1$ queries.
Then, $\sum_{x_t\in A,\zt=\zhf}\sigmasqtmo(x_t)\leq \frac{2}{\log(1+\eta^{-2})}\IGnn{s}(A)$.
Here $\IGnn{s}(A)$ is the MIG of $\phix$ after $s$ queries to $A$ as given
in Definition~\ref{def:infGain}.
% \newcommand{\xAss}[1]{x_{(#1)}}
\end{lemma}

We now define the quantity $R_n$ below.
Readers familiar with the GP bandit literature might see that it is similar to the notion
of cumulative regret, but we only consider queries at $\zhf$.
\begin{align*}
R_n = \sum_{\substack{t=1\\ \zt=\zhf}}^{n} \funcopt - \func(\xt) \;\;=\;\;
  \sum_{\substack{\zt=\zhf\\\xt\in\Xcalrhon}}
        \funcopt - \func(\xt)
\quad + \quad
  \sum_{\substack{\zt=\zhf\\\xt\notin\Xcalrhon}}
        \funcopt - \func(\xt).
\numberthis
\label{eqn:Rn}
\end{align*}
% We first work with the first term on the right hand side.
% Using a 
For any $A\subset\Xcal$ we can use Lemmas~\ref{lem:ucblemma},~\ref{lem:smoothnesslemma},
and~\ref{lem:IGBoundLemma} and the Cauchy Schwartz inequality to obtain,
\[
\sum_{\substack{\zt=\zhf\\\xt\in A}}
      \funcopt - \func(\xt)
\leq \sqrt{C_1 \Tn(\zhf,A)\betan \IG_{\Tn(\zhf,A)}(A)} + 
\sum_{\substack{\zt=\zhf\\\xt\in A}} \frac{1}{t^2}.
\]
For the first term in~\eqref{eqn:Rn}, we use the trivial bound
$\Tn(\zhf,\Xcalrhon) \leq n$.
For the second term we use the fact that
$\{\zhf\}\times(\Xcal\setminus\Xcalrhon)\subset\Fcaln$ and hence,
$\Tn(\zhf,\Xcal\setminus\Xcalrhon) \leq \Tn(\Fcaln) \leq n^\alpha$.
% Further noting that $ \Xcal\setminus\Xcalrhon \subset \Xcal$, we have
Noting that $A\subset B\implies \IGn(A)\leq\IGn(B)$, we have
$R_n \leq \sqrt{C_1n\betan\IGn(\Xcalrhon)} + \sqrt{C_1n^\alpha\betan\IGnalpha(\Xcal)} 
+ \pi^2/6$.
Now, using the fact that $N \leq 2\ncapital$ for large enough $N$ we have,
\[
R_N \leq  \sqrt{{2C_1\ncapital\beta_{2\ncapital} \IGnn{2\ncapital}(\Xcalrhon)} }
 \;+\;
  \sqrt{{2^\alpha C_1\ncapital^\alpha\beta_{2\ncapital}
\IGnn{2\ncapital^\alpha}(\Xcal)}} 
  \;+\; \frac{\pi^2}{6}.
\]
The theorem now follows from the fact that $S(\capital) \leq \frac{1}{N}R_N$ by definition
and that $N\geq \ncapital$.
The failure instances arise out of Lemmas~\ref{lem:ucblemma},~\ref{lem:smoothnesslemma} and
the bound on $\Tn(\Fcaln)$, the summation of whose probabilities are bounded by
$\delta$.
\hfill\BlackBox

\vspace{0.2in}

\begin{remark}[Construction of covering for the SE kernel]
\label{rem:construction}
\emph{
We demonstrate that such a construction is always possible using the SE kernel.
Using the inequality $e^{-x} \geq 1 - x$ for $x>0$ we have,
\[
1 - \phix^2(\Dx)\phiz^2(\Dz) <
\frac{\Dx^2}{\hx^2} + 
\frac{\Dz^2}{\hz^2} 
\]
% Using the inequality $e^{-x} \leq 1 - x + x^2/2$ for $x>0$ we have,
% \[
% 1 - \phix^2(\Dx)\phiz^2(\Dz) > 
% \frac{\Dx^2}{\hx^2} + 
% \frac{\Dz^2}{\hz^2} -
% \frac{\Dx^4}{2\hx^4} - 
% \frac{\Dz^4}{2\hz^4}
% \]
where $\Dz,\Dx$ will be the $L_2$ diameters of the balls in the covering.
% We will be letting $\Dz,\Dx$ shrink with $n$, so for sufficiently large $n$ the
% right hand side will be larger than 
% $ \frac{\Dx^2}{2\hx^2} + \frac{\Dz^2}{2\hz^2}$.
Now let $h = \min\{\hz, \hx\}$ and choose 
\[
\Dx = \Dz = \frac{h}{2}\frac{\inffunc(\sqrt{p})}{\betanh} \costratio^{q},
\]
via which we have
$1 - \phiz^2(z) \phix^2(x) < \frac{1}{2}\inffunc(\sqrt{p})^2\costratio^{2q}/\betan$
as stated in the proof.
Noting that $\betan\asymp\logn$, using standard results on covering numbers, we can
show that the size of this covering will be
$\logn^{\frac{d+p}{2}}/\costratio^{q(d+p)}$.
A similar argument is possible for Mat\'ern kernels, but the exponent on $\logn$ will be
worse.
}
\end{remark}

\begin{remark}[Choice of $q$ for SE kernel]
\label{rem:qchoice}
\emph{
From the arguments in our proof and Remark~\ref{rem:construction}, we have that
the number of plays in a set $S\subset(\Zcal\times\Xcal)$ is
$T(S)\leq \vol(S) \log(n)^{\frac{d+p+2}{2}}
\left(\frac{\cost(\zhf)}{\cost_{min}}\right)^{q(p+d+2)}$.
However, we chose to work work $\cost_{min}$ mostly to simplify the proof.
It is not hard to see that for $A\subset\Xcal$ and $B\subset\Zcal$
if $\cost(z) \approx \cost'$ for all $z\in B$, then
$\Tn(B,A)\approx \vol(B\times A) \log(n)^{\frac{d+p+2}{2}}
\left(\frac{\cost(\zhf)}{\cost'}\right)^{q(p+d+2)}$.
As the capital spent in this region is $\cost'\Tn(A,B)$,
by picking $q=1/(p+d+2)$ we ensure that the 
capital expended for a certain $A\subset\Xcal$ at all fidelities is roughly the same,
i.e. for any $A$, the capital density in fidelities $z$ such that $\cost(z) <
\cost(\thetat(A))$ will be roughly the same.
\citet{kandasamy16mfbandit} showed that doing so achieved a nearly minimax optimal strategy
for cumulative regret in $K$-armed bandits.
While it is not clear that this is the best strategy for optimisation under GP
assumptions, it did reasonably well in our experiments. We leave it to
future work to resolve this.
}
\end{remark}


\subsubsection{Proof of Lemma~\ref{lem:gtrthetatbound}}
\label{sec:gtrthetatbound}
For brevity, we will denote $\theta = \thetat(A)$.
We will invoke the discretisation $G_t$ used in Lemma~\ref{lem:smoothnesslemma} via which
we have $\utilt([\xopt]_t) \geq \funcopt - 1/t^2$ for all $t\geq 1$.
Let $b = \argmax_{x\in A} \utilt(x)$ be the maximiser of the upper
confidence bound $\utilt$ in $A$ at time $t$.
Now note that,
$\xt\in A \implies \utilt(b) > \utilt([\xopt]_t) \implies
\utilt(b) > \funcopt - 1/t^2$.
We therefore have,
\begingroup
\allowdisplaybreaks
\begin{align*}
\PP\big(\Tnzzxx{\greaterzz{\theta}}{A} \,>\, u \big)\;
&\leq\; 
% \frac{\delta}{\pi^2}\frac{1}{u^{1+4/\alpha}}
\PP\big( \exists t: u+1\leq t \leq n, \;\;
  \utilt(b) > \funcopt - 1/t^2 \;\;\wedge\;\; \tautmo(b) <\gamma(\theta)) 
\\
&\leq\;
  \sum_{t=u+1}^n
\PP\big( \mutmo(b) - \func(b) >  \funcopt - \func(b) 
  - \betath\sigmatmo(b) - 1/t^2
\;\;\wedge\;\; \tautmo(b) <\gamma(\theta)\big) 
\numberthis\label{eqn:sumprobgtr}
% &\hspace{0.2in}\leq\;
%   \sum_{t=u+1}^n \PP_{Z\sim\Ncal(0,1)}\left(Z > (\rho_0-1)\betath\right)
% \;\leq\;
%   \sum_{t=u+1}^n \frac{1}{2}\exp\left(\frac{(\rho_0-1)^2}{2}\betat\right) \\
% &\hspace{0.2in}\leq\;
%   \frac{1}{2}\left(\frac{\delta}{M\pi^2}\right)^{(\rho_0-1)^2}
%   \sum_{t=u+1}^n t^{-(\rho_0-1)^2(2+2d)} 
% \;\leq\;
% \frac{\delta}{M\pi^2} u^{-(\rho_0-1)^2(2+2d) + 1}
% \;\leq\;
% \frac{\delta}{\pi^2} \frac{1}{u^{1+4/\alpha}}.
\end{align*}
\endgroup
We now note that 
\[
\tautmo(b)<\gamma(\theta)\implies \sigmatmo(b) < \gamma(\theta) +
\sqrtks\inffunc(\theta) \leq 2\sqrtks\inffunc(\theta)
\leq \frac{1}{\betath\rho}(\funcopt - \func(b)).
\]
The first step uses Lemma~\ref{lem:inffuncbound}.
The second step uses the fact that $\gamma(\theta) = \sqrtks\inffunc(\theta)
(\cost(z)/\cost(\zhf))^{1/(p+d+2)} \leq \sqrtks\inffunc(\theta)$ and
the last step uses the definition of $\Qt(A)$ in~\eqref{eqn:Qt} whereby we have
$\funcopt - \func(x) \geq 2\rho\betath\sqrtks\inffunc(\theta)$.
Now we plugging this back into~\eqref{eqn:sumprobgtr}, we can bound each term in
the summation by,
\begin{align*}
&\PP\big( \mutmo(b) - \func(b) > (\rho-1)\betath\sigmatmo(b) - 1/t^2 \big) 
\;\leq\;
  \PP_{Z\sim\Ncal(0,1)}\left(Z > (\rho_0-1)\betath\right) \\
&\hspace{0.2in}\leq\;
  \frac{1}{2}\exp\left(\frac{(\rho_0-1)^2}{2}\betat\right) 
\;\leq\;
  \frac{1}{2}\left(\frac{2\delta}{\pi^2}\right)^{(\rho_0-1)^2}
  t^{-(\rho_0-1)^2(2+2d)} 
\;\leq\;
  \frac{\delta}{\pi^2}   t^{-(\rho_0-1)^2(2+2d)} .
\numberthis\label{eqn:termprobgtr}
% \;\leq\;
% \frac{\delta}{M\pi^2} u^{-(\rho_0-1)^2(2+2d) + 1}
% \;\leq\;
% \frac{\delta}{\pi^2} \frac{1}{u^{1+4/\alpha}}.
\end{align*}
In the first step we have used the following facts,
$t>u\geq\max\{3,\,(2(\rho-\rho_0)\eta)^{-2/3}\},\,$
$\pi^2/2\delta >1$ and
$\sigmatmo(b) > \eta/\sqrt{t}$ to conclude,
\begin{align*}
(\rho-\rho_0)\frac{\eta\sqrt{4\logt}}{\sqrt{t}} > \frac{1}{t^2}
\;&\implies\;
(\rho-\rho_0)\cdot
\sqrt{2\log\left(\frac{\pi^2t^2}{2\delta}\right)}\cdot\frac{\eta}{\sqrt{t}} > 
            \frac{1}{t^2} 
\;\implies\;
(\rho-\rho_0) \betath\sigmatmo(b) > \frac{1}{t^2}.
\end{align*}
The second step of~\eqref{eqn:termprobgtr} uses Lemma~\ref{lem:gaussConcentration},
the third step uses the conditions on $\betath$ as given in theorem~\ref{thm:main}
and the last step uses the fact that $\pi^2/2\delta > 1$.
Now plug~\eqref{eqn:termprobgtr} back into~\eqref{eqn:sumprobgtr}.
The result follows by bounding the sum by an integral and  noting that
$\rho_0>2$ and $\rho_0\geq 1+ \sqrt{(1+2/\alpha)/(1+d)}$.
\hfill\BlackBox


\subsubsection{Proof of Lemma~\ref{lem:IGBoundLemma}}

\newcommand{\xAss}[1]{u_{#1}}
\newcommand{\xAt}{\xAss{t}}
\newcommand{\sigmattmo}{\tilde{\sigma}_{t-1}}
\newcommand{\sigmattmosq}{\tilde{\sigma}^2_{t-1}}
Let $A_s = \{\xAss{1},\xAss{2},\dots,\xAss{s}\}$ be the queries in $\{\zhf\}\times A$
in the order they were queried.
Now, assuming that we have queried $\gunc$ only inside $\{\zhf\}\times A$, denote by
$\sigmattmo(\cdot)$, the posterior standard deviation after $t-1$ such queries. Then,
\begin{align*}
\sum_{t:x_t\in A, \zt=\zhf} \sigmasqtmo(x_t) \leq\;
  \sum_{t=1}^s\sigmattmosq(\xAt) 
  \leq\; \sum_{t=1}^s \eta^2\frac{\sigmattmosq(\xAt)}{\eta^2}
  \leq\; \sum_{t=1}^s
    \frac{\log(1+\eta^{-2}\sigmattmosq(\xAt))}{\log(1+\eta^{-2})} 
  \leq \frac{2}{\log(1+\eta^{-2})} I(y_{A_s}; f_{A_s}).
\end{align*}
Queries outside $\{\zhf\}\times A$ will only decrease the variance of the GP so we can
upper bound
the first sum by the posterior variances of the GP with only the queries in
$\{\zhf\}\times A$.
The third step uses the inequality $u^2/v^2 \leq \log(1+u^2)/\log(1+v^2)$.
The result follows from the fact that $\IGnn{s}(A)$
maximises the mutual information among all subsets of size $s$.
\hfill\BlackBox


